% Generated draft: requires mathematical and source review.
\documentclass[11pt]{article}
\usepackage[margin=25mm]{geometry}
\usepackage{fontspec}
\setmainfont{Latin Modern Roman}
\usepackage{amsmath,amssymb,mathtools,unicode-math}
\setmathfont{Latin Modern Math}
\usepackage{enumitem,needspace,xurl,hyperref}
\setlength{\parindent}{0pt}
\setlength{\parskip}{5pt}
\setlength{\emergencystretch}{3em}
\newcommand{\N}{\mathbb{N}}
\newcommand{\Z}{\mathbb{Z}}
\newcommand{\Q}{\mathbb{Q}}
\newcommand{\R}{\mathbb{R}}
\newcommand{\C}{\mathbb{C}}
\newcommand{\F}{\mathbb{F}}
\newcommand{\A}{\mathbb{A}}
\newcommand{\PP}{\mathbb{P}}
\newcommand{\E}{\mathbb{E}}
\newcommand{\eps}{\varepsilon}
\newcommand{\dd}{\,\mathrm{d}}
\begin{document}
\section*{UnsolvedMath ID: 6700077}
\textbf{Source problem code:} catalogue entry
\subsection{Definitions and mathematical statement}
Let \(\operatorname{Met}^0(X)\) be the continuous positive-definite symmetric two-tensor fields on a smooth \(n\)-manifold \(X\), with \(C^0\) convergence interpreted as uniform convergence on compact subsets in any fixed smooth background norm. For the source's volumic scalar-curvature notion, set provisionally
\[
\mathcal C^{\mathrm{vol}}_\kappa(X)=
\{g\in\operatorname{Met}^0(X):\operatorname{Sc}^{\mathrm{vol}}_n(g)\ge\kappa\}.
\]
The inequality is the intended lower bound at every point. The definition and normalization of \(\operatorname{Sc}^{\mathrm{vol}}_n\), including its admissible \(\kappa\)-range, are missing from the extraction.

The conjecture is closure, not density:
\[
g_j\in\mathcal C^{\mathrm{vol}}_\kappa(X)\ (j\ge1),\qquad
 g_j\longrightarrow g\text{ in }C^0
\quad\Longrightarrow\quad
 g\in\mathcal C^{\mathrm{vol}}_\kappa(X).
\]
The same fixed lower bound \(\kappa\) is retained in the limit; the approximating metrics and the limit are only required to be continuous. Equivalently, the indicated subset of continuous metrics is closed in the chosen topology.

The source motivates this with a reported smooth closure theorem: if \(g,g_j\) are \(C^2\) metrics, \(g_j\to g\) in \(C^0\), and \(\operatorname{Sc}_{g_j}\ge\kappa\), then \(\operatorname{Sc}_g\ge\kappa\), citing [49,10]. This uses ordinary smooth scalar curvature and does not itself define the continuous volumic predicate.

The supplied research summary instead discusses flexibility and possible density of volumically positive metrics, sometimes only almost everywhere. Such results or conjectures are logically different from the displayed closure statement and from a pointwise lower bound; they are not substituted for it. The stored text ends with an unfinished Test. Check [? cross-reference, whose missing example or test cannot be reconstructed.
\subsection{Short English statement}
Is a fixed volumic scalar-curvature lower bound preserved under \(C^0\) limits of continuous Riemannian metrics? This is a closure question; the supplied density/flexibility discussion does not answer it.
\subsection{Sources}
\begin{itemize}
\item Local UnsolvedMath registry, permanent ID \(6700077\), exact source code \texttt{catalogue entry}; printed scalar-curvature tag [?79], occurrence 80. Proposer Misha Gromov, 2017.
\item Gromov, \emph{101 Questions, Problems and Conjectures around Scalar Curvature} (2017), recorded PDF page 81; \url{https://www.ihes.fr/~gromov/wp-content/uploads/2018/08/101-problemsOct1-2017.pdf}. Recorded access 2026-07-29; pdf-text-normalized extraction.
\item Source references [49,10] for the ordinary \(C^2\)-metric scalar-curvature closure theorem; bibliography entries are absent. The final test reference is incomplete.
\item Supplied research summary mentions Gromov's volumic flexibility constructions but provides no precise citation or theorem establishing the stated continuous-metric closure conjecture.
\end{itemize}
\end{document}
